\chapter{TT}
% \section{$\lambda D$ System}
% 
% As we did in FOL, we first need to introduce some symbols in order to construct a formal language of TT.
% 
% \begin{itemize}
%   \item $V = \{ v_0, v_1, \cdots \}$ is the set of term variables;
%   \item $A = \{ \alpha_0, \alpha_1, \cdots \}$ is the set of type variables (below we take the union of $A$ and $V$, and denote it by $V$);
%   \item $C = \{ c_0, c_1, \cdots \}$ is the set of constants that are new definitions;
%   \item $P = \{ ., :, \equiv, (, ) \}$ is the set of punctuation;
%   \item $B = \{ \lambda, \Pi \}$ is the set of binders;
%   \item $H = \{ *, \square \}$ is the set of higher order types.
% \end{itemize}
% 
% Then we let
% \[
%   \Fml = V | H | \Fml\Fml | B V: \Fml.\Fml | C\Fml^*
% \]
% in which $\Fml^*$ is the set of words over set $\Fml$, including the empty word $\emty$.
% 
% \begin{definition}
%   A definition $D$ has the form
%   \[
%     \begin{array}{ll}
%       \bm{x} : \bm{A} & c\bm{x} \equiv d(\bm x) : a 
%     \end{array}
%   \]
%   where we denote $\bm x $ by $ (x_i)_{i=0}^s$ and $\bm x : \bm A$ by 
%   \[
%     \begin{array}{llll}
%       x_0 : A_0 & x_1 : A_1 &\cdots &x_s : A_s
%     \end{array}
%   \]
%   note that $d(\bm x)$ refers $d \in \Fml$ and $\Fr d \subset \bm x$, that $x_i \in V, c \in C$ and that all $A_i, d, a$ are in $\Fml$.
% 
%   A environment $\Delta$ has the form $D_0 D_1 \cdots D_n$.
% \end{definition}
% 
% 
% 
% 
% 
% \begin{definition}
%   Given $\Delta = D_0 D_1 \cdots D_n$, $R_{\Delta} \subset \Fml^2$ is the smallest relation satisfies
%   the following properties:
%   \begin{itemize}
%     \item contains $\{ (c \bm U, d(\bm U)): \bm x : \bm A \ c \bm x \equiv d(\bm x) : a \in D \in \Delta \textrm{ and } \bm U \in \Fml^* \}$;
%     \item closed under the following operations: $\Fml -, - \Fml, \lambda V: \Fml. -, C\bm p_1 - \bm p_2$.
%   \end{itemize}
% \end{definition}
% 
% 
% 
% 
% \begin{definition}
%   $R_{\beta}$ is the smallest relation satisfies
%   \begin{itemize}
%     \item contains $\{ ((\lambda x: a . \varphi)\chi, \varphi \frac{\chi}{x}) : x \in V \textrm{ and all } \varphi, \chi, a \in \Fml \}$;
%     \item closed under the following operations: $\Fml -, - \Fml, B V: \Fml.-, B V: -. \Fml, C \bm p_1 - \bm p_2$.
%   \end{itemize}
% \end{definition}
% 
% 
% 
% 
% 
% \begin{definition}
%   $\sim_{\Delta, \beta} \subset \Fml^2$ is the equivalence generated by $R_{\Delta} \cup R_{\beta}$.
% \end{definition}
% 
% 
% 
% 
% 
% \begin{definition}
%   A judgment has the form
%   \[
%     \begin{array}{lll}
%       \Delta & \bm x : \bm A & \varphi : a
%     \end{array}
%   \] 
% \end{definition}
% 
% 
% 
% 
% All legal judgments are derived through a series of rules, as listed in follows. By the way, the question that whether a judgment is legal is decidable. Below we let $s \in H$ if it is not specified. 
% 
% 
% \subsubsection*{(Sort)}
% \[
%   \begin{array}{lll}
%     \emty & \emty & * : \square
%   \end{array}
% \]
% 
% 
% \subsubsection*{(Var)}
% If $x \in V \smallsetminus \bm x$, then
% \[
%   \begin{array}{llll}
%     \Delta & \bm x : \bm A & a : s & 
%     \\
%     \Delta & \bm x : \bm A & x : a & x : a
%   \end{array}
% \]
% 
% 
% \subsubsection*{(Weak)}
% If $x \in V \smallsetminus \bm x$, then
% \[
%   \begin{array}{llll}
%     \Delta & \bm x : \bm A & \varphi : a & 
%     \\ 
%     \Delta & \bm x : \bm A & b : s  &
%     \\
%     \Delta & \bm x : \bm A & x : b & \varphi : a
%   \end{array}
% \]
% 
% \subsubsection*{(Form)}
% \[
%   \begin{array}{llll}
%     \Delta & \bm x : \bm A & a_1 : s_1 & 
%     \\ 
%     \Delta & \bm x : \bm A & x : a_1   &  a_2 : s_2
%     \\
%     \Delta & \bm x : \bm A & \Pi x : a_1. a_2 : s_2 
%   \end{array}
% \]
% 
% 
% 
% 
% \subsubsection*{(Appl)}
% \[
%   \begin{array}{llll}
%     \Delta & \bm x : \bm A & \varphi : \Pi x : a . b
%     \\ 
%     \Delta & \bm x : \bm A & \chi : a 
%     \\
%     \Delta & \bm x : \bm A & \varphi \chi : b\frac{\chi}{x} 
%   \end{array}
% \]
% 
% 
% \subsubsection*{(Abst)}
% \[
%   \begin{array}{llll}
%     \Delta & \bm x : \bm A & x : a & \varphi : b
%     \\ 
%     \Delta & \bm x : \bm A & \Pi x : a.b : s &
%     \\
%     \Delta & \bm x : \bm A & (\lambda x : a. \varphi) : (\Pi x: a.b) & 
%   \end{array}
% \]
% 
% 
% \subsubsection*{(Def)}
% \[
%   \begin{array}{lll}
%     \Delta & \bm x : \bm A & \varphi : a 
%     \\ 
%     \Delta & \bm y : \bm B & d(\bm y) : b
%     \\
%     \Delta D & \bm x : \bm A  & \varphi : a 
%   \end{array}
% \]
% where $D$ is the string
% \[
%   \begin{array}{ll}
%     \bm y : \bm B & c \bm y \equiv d(\bm y) : b
%   \end{array}
% \]
% that $d(\bm y)$ could be an empty word.
% 
% 
% \subsubsection*{(Inst)}
% If $\Delta$ contains the element
% \[
%   \begin{array}{ll}
%     \bm y : \bm B & c \bm y \equiv d(\bm y) : a
%   \end{array}
% \]
% that $d(\bm y)$ could be empty, and we have $\bm U \in \Fml^*$ with the same size of $\bm y$, then 
% \[
%   \begin{array}{lll}
%     \Delta & \bm x : \bm A & * : \square 
%     \\ 
%     \Delta & \bm x : \bm A & U_0 : B_0 \frac{\bm U}{\bm y}
%     \\
%            &  \vdots       &
%     \\
%     \Delta & \bm x : \bm A & U_m : B_m \frac{\bm U}{\bm y}
%     \\
%     \Delta & \bm x : \bm A  & c \bm U : a \frac{\bm U}{\bm y}
%   \end{array}
% \]
% 
% 
% 
% 
% \subsubsection*{(Conv)}
% If $a \sim_{\Delta, \beta} a'$, then
% \[
%   \begin{array}{llll}
%     \Delta & \bm x : \bm A & \varphi : a
%     \\ 
%     \Delta & \bm x : \bm A & a' : s
%     \\
%     \Delta & \bm x : \bm A & \varphi : a' 
%   \end{array}
% \]
% 
% 
% 
% 
% \subsubsection*{(Derived: Par)}
% Given $\Delta$ and let $D$ be 
% \[
%   \begin{array}{ll}
%     \bm y: \bm B & c \bm y \equiv d(\bm y) : a
%   \end{array}
% \]
% satisfies that $c$ does not occur in $\Delta$, then
% \[
%   \begin{array}{lll}
%     \Delta & \bm x : \bm A & d(\bm y) : a
%     \\ 
%     \Delta D & \bm x : \bm A & c \bm y : a 
%   \end{array}
% \]
% 









\section{CIC System}

We first introduce some symbols in order to construct a formal language:
\begin{itemize}
  \item $V = \{ v_0, v_1, \cdots \}$ is the set of term variables;
  \item $A = \{ \alpha_0, \alpha_1, \cdots \}$ is the set of type variables (below we take the union of $A$ and $V$, and denote it by $V$);
  \item $C = \{ \text{all UTF-8 strings} \} $ is the set of constants that are new definitions;
  \item $I = \{ \Ind, \Cstr, \Elim \}$ is the set of symbols used in defining inductions;
  \item $S = \{\Set, \Prop\} \cup \{ \Type_i \}_{i \geq 0} $ is the set of sorts;
  \item $P = \{ ., :, \equiv, (, ) \}$ is the set of punctuation;
  \item $B = \{ \lambda, \Pi \}$ is the set of binders.
\end{itemize}
Then we let
\[
  \Fml = C | V | S | (BV:\Fml.\Fml) | (I \Fml) | (\Fml \Fml)
\]


Axioms $\cl A \subset S^2$ and rules $\cl R \subset S^3$ are defined as follows
\begin{gather*}
  \cl A = \{ (\Set, \Type_0), (\Prop, \Type_0), (\Type_i, \Type_{i+1}) \}\\
  \cl R = \{ (\Prop, \Prop, \Prop), (\Set, \Prop, \Prop), (\Type_i, \Prop, \Prop), \\ 
  (\Prop, \Set, \Set), (\Set, \Set, \Set), (\Type_i, \Type_j, \Type_{\max \{i, j\}}) \}.
\end{gather*}



Note that $\Pi (\bm x : \bm A) \varphi$ is the abbreviation of $\Pi x_0 : A_0 . \Pi x_1 : A_1 \cdots \Pi x_m : A_m. \varphi$. $\bm x : \bm A  $ is the abbreviation of $x_0 : A_0 x_1 : A_1 \cdots x_m : A_m$. 





\begin{definition}
  The Arity has the form
  \[
    \mr{A} = S | (\Pi V : \Fml. \mr{A})
  \]

  For any $X \in C$, the Strictly Positive Types has the form
  \[
    \mr P_{X} = X | (\mr P_{X} \underbracket{\Fml}_{X \notin \Fml})  | (\Pi V : \underbracket{\Fml}_{X \notin \Fml} . \mr P_{X} )
  \] 
  The Constructors has the form
  \[
    \mr C _X  = X |( \mr C_X \underbracket{\Fml}_{X \notin \Fml}) |( \Pi V : \underbracket{\Fml}_{X \notin \Fml} . \mr C_X ) | ( \Pi V : \mr P_X . \underbracket{\mr C_X}_{V \notin \mr C_X})
  \]
  Any element $\mr C_X \smallsetminus \mr P_X$ is called recursive. 
\end{definition}






\begin{definition}
  A definition $D$ has the form
  \[
    \begin{array}{ll}
      \bm{x} : \bm{A} & c\bm{x} \equiv d(\bm x) : a 
    \end{array}
  \]
  note that $d(\bm x)$ refers $d \in \Fml$ and $\Fr d \subset \bm x$, that $x_i \in V, c \in C$ and that all $A_i, d, a$ are in $\Fml$.

  A environment $\Delta$ has the form $D_0 D_1 \cdots D_n$.
\end{definition}






\begin{definition}
  A judgment has the form
  \[
    \begin{array}{lll}
      \Delta & \bm x : \bm A & \varphi : a
    \end{array}
  \] 
\end{definition}







\begin{definition}
  Given $\Delta = D_0 D_1 \cdots D_n$, $R_{\Delta} \subset \Fml^2$ is the smallest relation satisfies
  the following properties:
  \begin{itemize}
    \item contains $\{ (c \bm U, d(\bm U)): \bm x : \bm A \ c \bm x \equiv d(\bm x) : a \in D \in \Delta \textrm{ and } \bm U \in \Fml^* \}$;
    \item closed under the following operations: $\Fml -, - \Fml, \lambda V: \Fml. -, C\bm p_1 - \bm p_2$.
  \end{itemize}
\end{definition}




\begin{definition}
  $R_{\beta}$ is the smallest relation satisfies
  \begin{itemize}
    \item contains $\{ ((\lambda x: a . \varphi)\chi, \varphi \frac{\chi}{x}) : x \in V \textrm{ and all } \varphi, \chi, a \in \Fml \}$;
    \item closed under the following operations: $\Fml -, - \Fml, B V: \Fml.-, B V: -. \Fml, C \bm p_1 - \bm p_2$.
  \end{itemize}
\end{definition}




\begin{proposition}
  Any element in $\mr C_X$ and $\mr P_X$ has the form 
  \[
    \Pi (\bm x : \bm M) . (X \bm N)
  \]
  which is unique \textbf{under} $\sim_{\beta, \Delta}$.
\end{proposition}








Suppose $X$ is a constant symbol, $c_i \in \mr C_X$, and $\Pi (\bm x: \bm A).s \in \rm A$, then
\[
  \begin{array}{llll}
    \Delta   &\Gamma   &X: \Pi (\bm x: \bm A).s   &c_0 : s\\
             &         &\vdots                    &       \\
    \Delta   &\Gamma   &X: \Pi (\bm x: \bm A).s   &c_m : s\\
    \Delta   &\Gamma   &\Ind X :  \Pi (\bm x: \bm A).s   &
  \end{array}
\]
and it immediately follows that
\[
  \begin{array}{lll}
    \Delta   &\Gamma   &\Ind X :  \Pi (\bm x: \bm A).s \\
    \Delta   &\Gamma   &\Cstr \Ind X i : c_i \frac{\Ind X}{X}
  \end{array}
\]
($i$ is a constant belongs to $C$)





% Assume that $ X : \Pi (\bm x: \bm A).s $, $Y :  \Pi (\bm x: \bm A).s'$ and $c \in \mr C_X$. 

% This induces a function $E_{X, Y}: \mr C_X \to \mr C_Y$, which is defined inductively.
% \begin{align*}
%   &E_{X, Y}(X) = X\\
%   &E_{X, Y}(c M) =  E_{X, Y}(c) M\\
%   &E_{X, Y}(M \to c) = M \to M[Y/X] \to E_{X, Y}(c) \\
%   &E_{X, Y}(\Pi x: M.c) = \Pi x: M. E_{X, Y}(c) \ (X \notin M)
% \end{align*}
% 
% 
% \begin{lemma}
%   If for all $s''$, $(s'', s) \in \cl A$ implies $(s'', s') \in \cl A$, and $(s', s') \in \cl A$, then the follows is derivable
%   \[
%     \begin{array}{lllll}
%       \Delta   &\Gamma   &X :  \Pi (\bm x: \bm A).s   & c:s  & \\
%       \Delta   &\Gamma   &\Pi (\bm x: \bm A).s' & &\\
%       \Delta   &\Gamma   &X :  \Pi (\bm x: \bm A).s   &Y :  \Pi (\bm x: \bm A).s'  & E_{X, Y} (c) : s'
%     \end{array}
%   \]
% \end{lemma}
% 
% 
% 
% If $f_0, \cdots, f_m$ and $G$ are new constant symbols, we have
% \[
%   \begin{array}{lll}
%     \Delta   &\Gamma   & G: \Ind X \bm a  \\
%     \Delta   &\Gamma   & Y : \Pi (\bm x: \bm A).s' \\
%     \Delta   &\Gamma   & f_0 : E_{X, Y}(c_0) \\
%              &         & \vdots \\
%     \Delta   &\Gamma   & f_m : E_{X, Y}(c_m) \\
%     \Delta   &\Gamma   & \Elim G Y f_0 f_1 \cdots f_m : Y \bm a
%   \end{array}
% \]
% 



% Assume that $ X : \Pi (\bm x: \bm A).s $, $Y :  \Pi (\bm x: \bm A).s'$, $d : c \in \mr C_X$ and $p : P = \Pi(\bm x' : \bm P') . (X \bm m') \in \mr P_X$, we define a new function $E_{X, Y, d} : \mr C_X \to \mr C_Y$, as follows
Assume that 

\begin{itemize}
  \item $\Fml \ni X : \Pi (\bm x: \bm A).s $
  \item $C \ni Y :  \Pi (\bm x: \bm A).(X \bm x) \to s'$
  \item $\Fml \ni d : c \in \mr C_X$
  \item $P = \Pi(\bm x' : \bm P') . (X \bm m') \in \mr P_X$
  \item $p \in V$ is a new variable
\end{itemize}
we can define a new function $E_{X, Y, d} : \mr C_X \to \mr C_Y$ inductivly, as follows



\begin{align*}
  & E_{X, Y, d}(X) = Y \\ 
  & E_{X, Y, d}(c M) = (E_{X, Y, d}(c) M) d \\
  & E_{X, Y, d}(P \to c) = \Pi p : P . (\Pi (\bm x' : \bm P') . Y \bm m' (p \bm x')  \to E_{X, Y, d p} (c)) \\
  & E_{X, Y, d}(\Pi x : M. c) = \Pi x : M . E_{X, Y, d x}(c)
\end{align*}








% \begin{lemma}
%   If for all $s''$, $(s'', s) \in \cl A$ implies $(s'', s') \in \cl A$, and $(s', s') \in \cl A$, then the follows is derivable
%   \[
%     \begin{array}{lllll}
%       \Delta   &\Gamma   &X :  \Pi (\bm x: \bm A).s   & c:s  & \\
%       \Delta   &\Gamma   &(\Pi (\bm x: \bm A).(X \bm x)) \to s' & &\\
%       \Delta   &\Gamma   &X :  \Pi (\bm x: \bm A).s   &Y :  (\Pi (\bm x: \bm A).(X \bm x)) \to s'  & E_{X, Y, d} (c) : s'
%     \end{array}
%   \]
% \end{lemma}
% 
% 
% 






If $f_0, \cdots, f_m$ and $G$ are new constant symbols, we have
\[
  \begin{array}{lll}
    \Delta   &\Gamma   & G: \Ind X \bm a  \\
    \Delta   &\Gamma   & Y : \Pi (\bm x: \bm A) . ((\Ind X \bm x) \to s') \\
    \Delta   &\Gamma   & f_0 : E_{\Ind X, Y, \Cstr \Ind X 0}(c_0) \\
             &         & \vdots \\
    \Delta   &\Gamma   & f_m : E_{\Ind X, Y, \Cstr \Ind X m}(c_m) \\
    \Delta   &\Gamma   & \Elim G Y f_0 f_1 \cdots f_m : (Y \bm a G)
  \end{array}
\]







Again, suppose $u, v \in \Fml$ we define a function $H_{X, u, v}: \mr C_X \to \Fml$, by letting
\begin{align*}
  & H_{X, u ,v} (X) = v \\
  & H_{X, u ,v} (c M) = H_{X, u ,v} (c) \\
  & H_{X, u ,v} (P \to c) = \lambda p: P . H_{X, u ,(v p \lambda (\bm x : \bm P). (v \bm m (p \bm x)))}(c) \\
  & H_{X, u ,v} (\Pi x : M. c) = \lambda x : M . H_{X, u ,v x} (c)
\end{align*}




\begin{definition}
  Suppose $\bm f$ has the shape of $m$, $\rightarrow_{\iota}$ is a relation in $\Fml^2$ contains
  \[
  \Elim \Cstr \Ind X i \bm m Y \bm f \rightarrow_{\iota} H_{\Ind X, \lambda(\bm x : \bm A). \lambda G : \Ind X \bm x . \Elim G Y \bm f , f_i } (c_i)\bm m
  \]
\end{definition}








In COC, we have introduced the following sequent rules.



\subsubsection*{(Sort)}
If $(s_1, s_2) \in \cl A$, then 
\[
  \begin{array}{lll}
    \emty & \emty & s_1 ; s_2
  \end{array}
\]


\subsubsection*{(Var)}
If $x \in V \smallsetminus \bm x$, then
\[
  \begin{array}{llll}
    \Delta & \Gamma & a : s & 
    \\
    \Delta & \Gamma & x : a & x : a
  \end{array}
\]


\subsubsection*{(Weak)}
If $x \in V \smallsetminus \bm x$, then
\[
  \begin{array}{llll}
    \Delta & \Gamma & \varphi : a & 
    \\ 
    \Delta & \Gamma & b : s  &
    \\
    \Delta & \Gamma & x : b & \varphi : a
  \end{array}
\]

\subsubsection*{(Form)}
If $(s_1, s_2, s_3) \in \cl R$, then 
\[
  \begin{array}{llll}
    \Delta & \Gamma & a_1 : s_1 & 
    \\ 
    \Delta & \Gamma & x : a_1   &  a_2 : s_2
    \\
    \Delta & \Gamma & \Pi x : a_1. a_2 : s_3 
  \end{array}
\]




\subsubsection*{(Appl)}
\[
  \begin{array}{llll}
    \Delta & \Gamma & \varphi : \Pi x : a . b
    \\ 
    \Delta & \Gamma & \chi : a 
    \\
    \Delta & \Gamma & \varphi \chi : b\frac{\chi}{x} 
  \end{array}
\]


\subsubsection*{(Abst)}
\[
  \begin{array}{llll}
    \Delta & \Gamma & x : a & \varphi : b
    \\ 
    \Delta & \Gamma & \Pi x : a.b : s &
    \\
    \Delta & \Gamma & (\lambda x : a. \varphi) : (\Pi x: a.b) & 
  \end{array}
\]


\subsubsection*{(Def)}
\[
  \begin{array}{lll}
    \Delta & \Gamma & \varphi : a 
    \\ 
    \Delta & \bm y : \bm B & d(\bm y) : b
    \\
    \Delta D & \Gamma  & \varphi : a 
  \end{array}
\]
where $D$ is the string
\[
  \begin{array}{ll}
    \bm y : \bm B & c \bm y \equiv d(\bm y) : b
  \end{array}
\]
that $d(\bm y)$ could be an empty word.


\subsubsection*{(Inst)}
If $\Delta$ contains the element
\[
  \begin{array}{ll}
    \bm y : \bm B & c \bm y \equiv d(\bm y) : a
  \end{array}
\]
that $d(\bm y)$ could be empty, and we have $\bm U \in \Fml^*$ with the same size of $\bm y$, then 
\[
  \begin{array}{lll}
    \Delta & \Gamma & s_1 : s_2  \ (\in \cl A) 
    \\ 
    \Delta & \Gamma & U_0 : B_0 \frac{\bm U}{\bm y}
    \\
           &  \vdots       &
    \\
    \Delta & \Gamma & U_m : B_m \frac{\bm U}{\bm y}
    \\
    \Delta & \Gamma  & c \bm U : a \frac{\bm U}{\bm y}
  \end{array}
\]





\subsubsection*{(Conv)}
If $a \sim_{\Delta, \beta} a'$, then
\[
  \begin{array}{llll}
    \Delta & \Gamma & \varphi : a
    \\ 
    \Delta & \Gamma & a' : s
    \\
    \Delta & \Gamma & \varphi : a' 
  \end{array}
\]






